\section{A local obstruction to a semidefinite lift}
\label{sec:obstruction}

We prove that the quadratic matrix map constructed in
Section~\ref{sec:separation} has an epigraph set $E_Q$ with no finite
semidefinite lift.  The proof has two parts.  First, a polynomial patch in
$E_Q$ admits a positive moment functional whose application at small scales
violates the defining inequality of $E_Q$ to order four.  Second, any finite
lift would force the resulting points to lie much closer to $E_Q$ than this
violation allows.  The second part is a local sum-of-squares obstruction,
in the sense of the methods developed by Scheiderer
\cite[Section~4]{Scheiderer2018}; here we give a direct proof using finite
Taylor polynomials of Gram factors.

We retain the spaces $H=L\oplus P$, $R=K^\perp\subset P$, and
$W=L\oplus R$, the positive semidefinite Hankel matrix $T$, and the data
$N,w,U,g$ of the preceding sections.  In particular,
\[
 \dim H=36,\qquad \dim L=8,\qquad \dim R=12,\qquad
 \rank T=20,\qquad r=20N-1.
\]
Write $J=\ker T$.  We use that $J\cap L=0$ and that the projection of $J$
to $P$ is $K$.  As in Section~\ref{sec:separation}, the symbol $U$ also
denotes a $36N$-by-$r$ matrix whose columns form a basis of the subspace
$U\subset H^N$.  By Lemma~\ref{lem:generic-wu}, orthogonal projection
induces an isomorphism
\begin{equation}
 \label{eq:obstruction-U-projection}
 U\longrightarrow w^\perp\cap W^N.
\end{equation}
Moreover, $w\in W^N$ has components spanning $W$, and
$w^Tgw<0$.  The admissible quadratic matrix map is
$Q(y)=U^TG(y)gG(y)U$, as defined in~\eqref{eq:constructed-Q}.

To keep track of repeated blocks, for an operator $S$ on $H$ we write
$S_N=I_N\otimes S$, an operator on $H^N$.  The same convention applies to
maps between subspaces of $H$.  We order the coordinates of $H^N$ as
$L^N\oplus P^N$ when using two-by-two block matrices.  Subscripts $LL$,
$LP$, and $PP$ on an operator on $H^N$ refer to these blocks.

\subsection{A polynomial patch and its moment deformation}

Let $v(s)\in H$, for $s\in\R^7$, be the column of homogeneous quadratic
monomials evaluated at $(1,s)$.  Thus its entries have affine degrees zero,
one, and two.  Let $y_s$ be evaluation at $(1,s)$ on homogeneous quartics,
and define
\begin{equation}
 \label{eq:patch-matrices}
 V(s)=I_N\otimes v(s)\in\R^{36N\times N},\qquad
 M(s)=U^TV(s)\in\R^{r\times N}.
\end{equation}
The moment identity gives $G(y_s)=V(s)V(s)^T$.  In the following lemma,
$E_Q$ is the set defined in Proposition~\ref{prop:epigraph}.

\begin{lemma}[Polynomial patch]
 \label{lem:polynomial-patch}
There is an open ball $\Omega_0\subset\R^7$ about zero such that $M(s)$
has rank $N$ for every $s\in\Omega_0$.  The polynomial map
\begin{equation}
 \label{eq:polynomial-patch}
 x(s)=\bigl(M(s)M(s)^T,\ \tr(gG(y_s)),\ y_s\bigr)
\end{equation}
has degree at most four and takes $\Omega_0$ into $E_Q$.  Its values satisfy
the trace inequality defining $E_Q$ with equality.
\end{lemma}

\begin{proof}
Decompose $w=w_L+w_P$ in $L^N\oplus P^N$.  Since the components of $w$
span $W$ and $R\ne0$, we have $w_P\ne0$; also $w_P\in R^N$.  Given
$\ell\in L^N$, the vector
\[
 \ell-\frac{\langle\ell,w_L\rangle}{\|w_P\|^2}w_P
\]
belongs to $w^\perp\cap W^N$.  By
\eqref{eq:obstruction-U-projection}, it is the projection to $W^N$ of
some vector of $U$.  Thus projection from $U$ to $L^N$ is surjective.
The columns of $V(0)$ form a basis for the $N$ constant-coordinate
directions in $L^N$.  If $M(0)a=0$, then $V(0)a$ is orthogonal to $U$;
surjectivity onto $L^N$ forces $V(0)a=0$, and hence $a=0$.  Therefore
$\rank M(0)=N$, and the same holds on a sufficiently small open ball.

Put $\Theta(s)=V(s)^TgV(s)$.  The definition of $Q$ gives
\[
 Q(y_s)=M(s)\Theta(s)M(s)^T.
\]
Consequently
\[
 \ran Q(y_s)\subset\ran M(s)=\ran(M(s)M(s)^T).
\]
Full column rank of $M(s)$ gives
\[
 M(s)^T(M(s)M(s)^T)^\dagger M(s)=I_N,
\]
and hence
\[
 \tr\bigl((M(s)M(s)^T)^\dagger Q(y_s)\bigr)
 =\tr\Theta(s)=\tr(gG(y_s)).
\]
This proves membership in $E_Q$.  Finally, $M(s)$ has degree at most two,
while $G(y_s)$ and $y_s$ have degree at most four, proving the degree
bound.
\end{proof}

For a linear functional $\Lambda$ on polynomials in seven variables through
degree four, its moment matrix is the symmetric matrix indexed by monomials
of degree at most two, with $(\alpha,\beta)$ entry
$\Lambda(s^{\alpha+\beta})$.  Under the affine chart $x_0=1$, this is
precisely a Hankel matrix on $H$.  Applying $\Lambda$ to a vector or matrix
of polynomials always means applying it entrywise.

\begin{lemma}[A deformation outside $E_Q$]
 \label{lem:moment-violation}
There exist a linear functional
$\Lambda:\R[s_1,\ldots,s_7]_{\le4}\to\R$ and an open ball
$\Omega\subset\Omega_0$ about zero with the following properties:
\begin{enumerate}
 \item $\Lambda(1)=1$, and its moment matrix on polynomials of degree at
 most two is positive definite;
 \item for every $c\in\Omega$, writing
 \begin{equation}
  \label{eq:tested-patch}
  (A_{c,\eps},t_{c,\eps},y_{c,\eps})
  =\Lambda_s\bigl(x(c+\eps s)\bigr),
 \end{equation}
 there are constants $a_c>0$ and $\kappa_c>0$ such that, as
 $\eps\downarrow0$,
 \begin{gather}
  \label{eq:deformation-lower-bound}
  A_{c,\eps}\succeq a_c\eps^4 I_r,\\
  \label{eq:deformation-violation}
  t_{c,\eps}-\tr\bigl(A_{c,\eps}^{-1}Q(y_{c,\eps})\bigr)
  =-\kappa_c\eps^4+o(\eps^4).
 \end{gather}
 All three entries of~\eqref{eq:tested-patch} remain bounded as
 $\eps\downarrow0$.
\end{enumerate}
\end{lemma}

\begin{proof}
We first allow the moment matrix to be an arbitrary Hankel matrix $S$
near $T$, and normalize its constant moment at the end.  Let $B_c$ be the
matrix on $H$ determined by
$B_cv(s)=v(s+c)$.  This is an invertible matrix, depends polynomially on
$c$, and satisfies $B_0=I_H$.  In the affine chart, $L$ consists of the
constant and linear monomials, and $P$ consists of the quadratic monomials.
On this degree decomposition of $H$,
put
\[
 D_\eps=\diag(1,\eps I_7,\eps^2I_{28}).
\]
The functional with moment matrix $S$, applied at $c+\eps s$, has moment
matrix $B_cD_\eps S D_\eps B_c^T$.  Its repeated matrix is therefore
\begin{equation}
 \label{eq:translated-scaled-moment}
 G=B_{c,N}D_{\eps,N}S_ND_{\eps,N}B_{c,N}^T.
\end{equation}
For the corresponding tested data, $A=U^TGU$ and $t=\tr(gG)$.
Whenever $A$ is invertible, the quantity to be made negative is
\begin{equation}
 \label{eq:trace-residual}
 t-\tr(A^{-1}Q(y))
 =\tr\bigl(g\,[G-GU(U^TGU)^{-1}U^TG]\bigr).
\end{equation}

We now analyze the matrix inside this trace at small $\eps$.  Projection
of the column space of $B_{c,N}^TU$ to $L^N$ remains surjective for $c$
near zero.  We can therefore choose a continuously varying basis matrix
$U_c$ for this column space in the form
\begin{equation}
 \label{eq:Uc-blocks}
 U_c=\begin{pmatrix}I_{8N}&0\\ F_c&E_c\end{pmatrix},
 \qquad
 F_c\in\R^{28N\times8N},\quad
 E_c\in\R^{28N\times(12N-1)}.
\end{equation}
For completeness, keep one invertible $8N$-by-$8N$ minor of the projection
matrix and normalize its columns to have $L^N$ components $I_{8N}$.
From each remaining column subtract its $L^N$ component expressed in
these normalized columns.  The resulting columns have zero $L^N$
component and give the block $E_c$.  All these operations depend
continuously on $c$.  At $c=0$, projection of the
columns of $E_0$ to $R^N$ is injective: a vector $(0,E_0z)$ with zero
$R^N$ projection would belong to $U\cap K^N$, which is zero by
\eqref{eq:obstruction-U-projection}.

Let $D_{\eps,L,N}$ denote the restriction of $D_{\eps,N}$ to $L^N$.
For $\eps>0$, multiplication of $D_{\eps,N}U_c$ on the right by the
invertible block diagonal matrix with blocks $D_{\eps,L,N}^{-1}$ and
$\eps^{-2}I_{12N-1}$ gives
\begin{equation}
 \label{eq:rescaled-U}
 \widehat U_{c,\eps}
 =\begin{pmatrix}
 I_{8N}&0\\
 \eps^2F_cD_{\eps,L,N}^{-1}&E_c
 \end{pmatrix}.
\end{equation}
This matrix extends continuously to $\eps=0$, where the lower-left block
is zero.  The Gram matrix
$\widehat U_{0,0}^TT_N\widehat U_{0,0}$ is positive definite.  Indeed,
if a vector $(\ell,E_0z)$ belongs to $J^N$, then $E_0z\in K^N$.
Injectivity of its $R^N$ projection gives $z=0$, and $J\cap L=0$ then
gives $\ell=0$.  Thus the column space of $\widehat U_{0,0}$ meets
$\ker T_N$ trivially.

It follows that, for $(c,\eps,S)$ near $(0,0,T)$, the Gram matrix in
\begin{equation}
 \label{eq:rescaled-residual}
 \widehat C_{c,\eps}(S)
 =S_N-S_N\widehat U_{c,\eps}
  \bigl(\widehat U_{c,\eps}^TS_N\widehat U_{c,\eps}\bigr)^{-1}
  \widehat U_{c,\eps}^TS_N
\end{equation}
is invertible.  In particular, the displayed residual is continuous in all
three arguments, including at $\eps=0$.  Changing a basis of the columns
does not change the residual in~\eqref{eq:trace-residual}.  After removing
the outer translation factors, define
\[
 C_{c,\eps}(S)=D_{\eps,N}\widehat C_{c,\eps}(S)D_{\eps,N}.
\]
The exact relation to the matrix in~\eqref{eq:trace-residual} is
\begin{equation}
 \label{eq:restored-residual}
 G-GU(U^TGU)^{-1}U^TG
 =B_{c,N}C_{c,\eps}(S)B_{c,N}^T.
\end{equation}
The matrix $C_{c,\eps}(S)$ annihilates $U_c$.  If a symmetric block matrix
$C$ annihilates the
first $8N$ columns $\binom{I_{8N}}{F_c}$ of $U_c$, its blocks satisfy
$C_{PL}=-C_{PP}F_c$ and $C_{LL}=F_c^TC_{PP}F_c$.  Hence, with
\[
 \iota_c=\begin{pmatrix}-F_c^T\\I_{28N}\end{pmatrix}
 :P^N\longrightarrow H^N,
\]
we obtain the exact identity
\begin{equation}
 \label{eq:order-four-residual}
 C_{c,\eps}(S)
 =\eps^4\iota_c\bigl(\widehat C_{c,\eps}(S)\bigr)_{PP}\iota_c^T
 \qquad(\eps>0).
\end{equation}
Combining~\eqref{eq:restored-residual} and
\eqref{eq:order-four-residual}, the limit of
\eqref{eq:trace-residual} divided by $\eps^4$ is
\begin{equation}
 \label{eq:continuous-leading-coefficient}
 \tr\bigl(gB_{c,N}\iota_c
       (\widehat C_{c,0}(S))_{PP}\iota_c^TB_{c,N}^T\bigr).
\end{equation}
This expression is continuous in $c$ and $S$.  Its sign can therefore be
computed first at the singular moment matrix $T$.

At $(c,\eps,S)=(0,0,T)$, the matrix in
\eqref{eq:rescaled-residual} is positive semidefinite of rank
$20N-r=1$.  To see this rank directly, set
$Z=T_N^{1/2}\widehat U_{0,0}$.  The residual equals
$T_N^{1/2}(I-P_Z)T_N^{1/2}$, where $P_Z$ is the orthogonal projection
onto the $r$-dimensional column space of $Z$, a subspace of
$\ran T_N^{1/2}$.  It annihilates $L^N$, the embedded columns
$\binom{0}{E_0}$, and
$J^N$.  Its range is consequently contained in
\[
 P^N\cap(J^N)^\perp\cap(\ran E_0)^\perp
 =R^N\cap(\ran E_0)^\perp.
\]
Because projection of $E_0$ to $R^N$ has rank $12N-1$, this last space is
one-dimensional.  It contains $w_P$, since $w$ is orthogonal to every
column of $U_0$.  Therefore
\[
 \bigl(\widehat C_{0,0}(T)\bigr)_{PP}
 =\alpha w_Pw_P^T\qquad\text{for some }\alpha>0.
\]
The other columns of $U_0$ give $w_L=-F_0^Tw_P$, so
$\iota_0w_P=w$.  Restoring the translation in
\eqref{eq:order-four-residual}, whose matrix is the identity at $c=0$,
shows that the limiting coefficient in~\eqref{eq:trace-residual} is
\[
 \alpha\tr(gww^T)=\alpha w^Tgw<0.
\]

Choose a positive definite Hankel matrix $S$ sufficiently close to $T$.
Such a choice is possible by adding a small positive multiple of the
degree-two moment matrix of a nondegenerate Gaussian measure on $\R^7$.
By the continuity just proved, the limiting coefficient remains negative
for every $c$ in an open ball about zero.  Its constant moment is positive.
Normalize the corresponding functional by its value at $1$ and call the
result $\Lambda$.  Scaling the functional by a positive number $\rho$
scales $A,t,y$ by $\rho$, and scales $Q(y)$ by $\rho^2$.  It therefore
scales~\eqref{eq:trace-residual} by $\rho$, preserving the sign.  Shrinking
the ball to lie in $\Omega_0$ proves~\eqref{eq:deformation-violation}.

Let $S_\Lambda\succ0$ be the normalized moment matrix.  For fixed $c$,
the matrix $B_{c,N}^TU$ has full column rank.  Since the smallest diagonal
entry of $D_{\eps,N}$ is $\eps^2$ for $0<\eps\le1$, formula
\eqref{eq:translated-scaled-moment} gives, for $z\in\R^r$,
\[
 z^TA_{c,\eps}z
 \ge \lambda_{\min}(S_\Lambda)\eps^4
       \|B_{c,N}^TUz\|^2
 \ge a_c\eps^4\|z\|^2
\]
with $a_c>0$.  Finally, the tested coordinates are polynomials in $\eps$,
so they remain bounded.  This completes the proof.
\end{proof}

The preceding lemma supplies points outside $E_Q$ whose violation has a
controlled leading term.  We next show that a semidefinite lift would
force any such moment deformation of a polynomial patch to be arbitrarily
well approximated, to a prescribed finite order, by feasible points.

\subsection{Moment tests of analytic lift data}

We first record why moments given only through degree four can be used to
test Taylor polynomials of higher degree.  The assertion concerns finite
positive moment matrices and does not require a representing measure.

\begin{lemma}[Finite positive extension]
 \label{lem:positive-moment-extension}
Let $n\ge1$ and $j\ge0$ be integers, let $z=(z_1,\ldots,z_n)$, and let the
linear functional
$\lambda:\R[z]_{\le2j}\to\R$ satisfy $\lambda(p^2)>0$ for every
nonzero $p\in\R[z]_{\le j}$.  For every integer $b\ge j$, there is an
extension $\widetilde\lambda:\R[z]_{\le2b}\to\R$ such that
$\widetilde\lambda(p^2)>0$ for every nonzero
$p\in\R[z]_{\le b}$.
\end{lemma}

\begin{proof}
It suffices to extend from $j$ to $j+1$.  Index the enlarged moment matrix
first by monomials of degree at most $j$, and then by monomials homogeneous
of degree $j+1$.  It has the form
\[
 \begin{pmatrix}M_j&C\\C^T&D\end{pmatrix},\qquad M_j\succ0.
\]
The mixed block $C$ uses moments of degree at most $2j+1$.  Its entries
of degree at most $2j$ are already fixed, and we may assign arbitrary
values to all moments of degree $2j+1$.  Every entry of $D$ has degree
exactly $2j+2$, so none of these moments has yet been assigned.  Give them
the values $a\int z^\gamma\,d\gamma_n(z)$, where $\gamma_n$ is a
nondegenerate Gaussian probability measure and $a>0$.  Then $D=aD_0$,
where $D_0\succ0$: the integral of the square of any nonzero homogeneous
polynomial is strictly positive.  Choosing $a$ sufficiently large makes
$aD_0-C^TM_j^{-1}C$ positive definite.  The Schur complement proves
positivity of the enlarged moment matrix.  Each degree layer is assigned
as a functional on monomials, so all repeated Hankel entries agree.
Iteration proves the assertion.
\end{proof}

The next lemma provides simultaneous analytic lift variables and Gram
factors.  Only a nonempty open part of the polynomial patch is needed.

\begin{lemma}[Local analytic Gram factors]
 \label{lem:analytic-lift}
Let $C\subset\R^a$ have an exact representation
\[
 C=\{x\in\R^a:\text{there exists }u\in\R^\nu
                         \text{ with }\mathcal L(x,u)\succeq0\},
\]
where $\nu\ge0$ and $\mathcal L$ is an affine symmetric pencil of size
$\ell\ge1$.
Let $\Omega\subset\R^n$ be nonempty, open, and semialgebraic, and let
$p:\R^n\to\R^a$ be polynomial with $p(\Omega)\subset C$.
There is a nonempty open subset $\Omega'\subset\Omega$ and real analytic
maps $u:\Omega'\to\R^\nu$ and
$F:\Omega'\to\R^{\ell\times\ell}$ such that
\begin{equation}
 \label{eq:analytic-Gram-identity}
 \mathcal L(p(s),u(s))=F(s)F(s)^T\qquad(s\in\Omega').
\end{equation}
\end{lemma}

\begin{proof}
Consider the semialgebraic set
\[
 \mathcal A=\{(s,u,F):s\in\Omega,\quad
                  \mathcal L(p(s),u)=FF^T\}.
\]
Its projection to $\Omega$ is surjective, since every positive
semidefinite matrix has a real square Gram factor.  A semialgebraic set
admits a finite stratification by real analytic semialgebraic manifolds;
see \cite[Proposition~9.1.8]{BCR1998}, or
\cite[Section~2.2.2]{LRT2025}.  On at least one stratum, the differential
of the projection has rank $n$ at some point.  Otherwise Sard's theorem
on each of the finitely many strata would imply that the whole image has
Lebesgue measure zero, contradicting that it contains $\Omega$.
The analytic submersion theorem at such a point gives a local analytic
section of the projection.  Its $u$ and $F$ coordinates give
\eqref{eq:analytic-Gram-identity}.
\end{proof}

\begin{lemma}[Feasible approximation after a moment test]
 \label{lem:moment-near-feasibility}
Let $C$, $p$, and $\Omega$ be as in Lemma~\ref{lem:analytic-lift}.
Fix an integer $b\ge\max\{1,\deg p\}$ and a linear functional
$\lambda:\R[z_1,\ldots,z_n]_{\le2b}\to\R$ satisfying
$\lambda(1)=1$ and $\lambda(q^2)\ge0$ for all
$q\in\R[z]_{\le b}$.  There is a nonempty open subset
$\Omega'\subset\Omega$ such that, for each $c\in\Omega'$, the points
\[
 x_\eps=\lambda_z\bigl(p(c+\eps z)\bigr)
\]
admit points $\widetilde x_\eps\in C$, for all sufficiently small
$\eps>0$, with
\begin{equation}
 \label{eq:near-feasibility}
 \|\widetilde x_\eps-x_\eps\|=O(\eps^b).
\end{equation}
The implied constant may depend on $c$ and on the displayed data.
\end{lemma}

\begin{proof}
Choose analytic maps $u,F$ on $\Omega'$ as in
Lemma~\ref{lem:analytic-lift}, and fix $c\in\Omega'$.  For a real analytic
map $f$, write its degree-$b$ Taylor polynomial at $c$, evaluated at
$c+\eps z$, as
\[
 f_{c,b}(\eps,z)=\sum_{|\alpha|\le b}
  \frac{\partial^\alpha f(c)}{\alpha!}\eps^{|\alpha|}z^\alpha.
\]
Set $u_\eps=\lambda_z(u_{c,b}(\eps,z))$ and
\[
 P_\eps=\mathcal L(x_\eps,u_\eps),\qquad
 H_\eps=\lambda_z\bigl(F_{c,b}(\eps,z)F_{c,b}(\eps,z)^T\bigr).
\]
For every $v\in\R^\ell$, the scalar $v^TH_\eps v$ is the value of
$\lambda$ on a sum of squares of polynomials of degree at most $b$.
Thus $H_\eps\succeq0$.  The Taylor identity
\eqref{eq:analytic-Gram-identity} implies that $P_\eps$ and $H_\eps$
have the same coefficients through degree $b$ in $\eps$.  Here
$\deg p\le b$ ensures that the polynomial patch is retained exactly,
and $\lambda(1)=1$ retains the constant term of the affine pencil.
Both tested expressions are polynomials in $\eps$, so
\begin{equation}
 \label{eq:approximate-pencil}
 \|P_\eps-H_\eps\|=O(\eps^{b+1}),\qquad
 P_\eps\succeq-O(\eps^{b+1})I_\ell.
\end{equation}

To turn this approximate positivity into exact feasibility, let
\[
 K_0=\bigcap_{\mathcal L(x,u)\succeq0}\ker\mathcal L(x,u),
 \qquad S_0=K_0^\perp.
\]
Every selected matrix $\mathcal L(p(s),u(s))$ annihilates $K_0$.
Its Taylor coefficients at $c$ therefore do so as well.  Since
$P_\eps$ is obtained by applying $\lambda$ to its Taylor polynomial
through degree $b$, we have
\begin{equation}
 \label{eq:exact-common-kernel}
 P_\eps K_0=0.
\end{equation}
In particular, any negative eigenvalues of $P_\eps$ occur on $S_0$.

There is a feasible pencil matrix $P_*=\mathcal L(x_*,u_*)$ whose
restriction to $S_0$ is positive definite.  Indeed, choose a feasible
matrix of maximum rank.  For any other feasible matrix $P$, the average
$(P_*+P)/2$ is feasible, and
$\ker(P_*+P)=\ker P_*\cap\ker P$.  Maximality of the rank therefore
forces $\ker P_*\subset\ker P$, proving $\ker P_*=K_0$.
If $S_0=0$, equation~\eqref{eq:exact-common-kernel} already gives
$P_\eps=0$, so $x_\eps\in C$ and there is nothing to prove.

Otherwise choose $\mu>0$ with $P_*|_{S_0}\succeq\mu I_{S_0}$ and put
$\delta_\eps=\eps^b$.  For sufficiently small positive $\eps$,
\[
 \bigl((1-\delta_\eps)P_\eps+\delta_\eps P_*\bigr)|_{S_0}
 \succeq\bigl(\mu\eps^b-C_1\eps^{b+1}\bigr)I_{S_0}
 \succeq0.
\]
Both summands annihilate $K_0$, so the mixed matrix is positive
semidefinite on all of $\R^\ell$.  By affineness of the pencil it is the
pencil value at
\[
 \widetilde x_\eps=(1-\delta_\eps)x_\eps+\delta_\eps x_*,\qquad
 \widetilde u_\eps=(1-\delta_\eps)u_\eps+\delta_\eps u_*.
\]
Thus $\widetilde x_\eps\in C$.  The vector $x_\eps$ is polynomial in
$\eps$ and hence bounded near zero, giving~\eqref{eq:near-feasibility}.
\end{proof}

\subsection{The contradiction}

\begin{theorem}
 \label{prop:no-lift}
For the admissible quadratic matrix map $Q$ of
\eqref{eq:constructed-Q}, the set $E_Q$ is not a spectrahedral shadow.
Consequently its associated entire closed hyperbolicity cone $K_Q$ is not
a spectrahedral shadow.
\end{theorem}

\begin{proof}
Suppose that $E_Q$ has a finite affine semidefinite lift.  Take the
polynomial patch $x$, the functional $\Lambda$, and the open ball
$\Omega$ from Lemmas~\ref{lem:polynomial-patch}
and~\ref{lem:moment-violation}.  By
Lemma~\ref{lem:positive-moment-extension}, extend $\Lambda$ through
degree $32$ so that it is positive definite on squares of polynomials of
degree at most $b=16$.  This extension leaves all tested patch coordinates
unchanged, since their degrees are at most four, and preserves
$\Lambda(1)=1$.

Apply Lemma~\ref{lem:moment-near-feasibility} with $C=E_Q$, $p=x$,
and $b=16$.  Fix a center $c$ in the resulting nonempty open subset of
$\Omega$, and abbreviate the tested data by
$x_\eps=(A_\eps,t_\eps,y_\eps)$.  There are feasible points
\[
 \widetilde x_\eps
 =(A_\eps+\Delta A_\eps,\ t_\eps+\Delta t_\eps,
                         y_\eps+\Delta y_\eps)\in E_Q
\]
with all three perturbations of norm $O(\eps^{16})$.
By~\eqref{eq:deformation-lower-bound},
$A_\eps\succeq a_c\eps^4 I_r$.  Therefore, for sufficiently small
$\eps>0$,
\[
 A_\eps+\Delta A_\eps\succeq\tfrac12a_c\eps^4 I_r,
 \qquad
 \|A_\eps^{-1}\|+
 \|(A_\eps+\Delta A_\eps)^{-1}\|=O(\eps^{-4}).
\]
The inverse identity gives
\begin{align*}
 &(A_\eps+\Delta A_\eps)^{-1}-A_\eps^{-1}\\
 &\hspace{1cm}=-(A_\eps+\Delta A_\eps)^{-1}
          \Delta A_\eps A_\eps^{-1}=O(\eps^8)
\end{align*}
in operator norm.  The vectors $y_\eps$ are bounded, and $Q$ is
quadratic, so
\[
 Q(y_\eps)=O(1),\qquad
 Q(y_\eps+\Delta y_\eps)-Q(y_\eps)=O(\eps^{16}).
\]
Combining these estimates and using that $r$ is fixed yields
\begin{align*}
 &t_\eps+\Delta t_\eps
  -\tr\bigl((A_\eps+\Delta A_\eps)^{-1}
                         Q(y_\eps+\Delta y_\eps)\bigr)\\
 &\qquad=t_\eps-\tr\bigl(A_\eps^{-1}Q(y_\eps)\bigr)
                +O(\eps^8)\\
 &\qquad=-\kappa_c\eps^4+o(\eps^4)<0
\end{align*}
for all sufficiently small positive $\eps$.  Since
$A_\eps+\Delta A_\eps$ is positive definite, this contradicts the
trace inequality defining the feasible point $\widetilde x_\eps\in E_Q$.
Thus $E_Q$ has no finite semidefinite lift.  Proposition~\ref{prop:epigraph}
identifies $E_Q$ as a projection of an affine section of $K_Q$, so the
same conclusion holds for $K_Q$.
\end{proof}

\begin{proof}[Proof of Theorem~\ref{thm:main}]
The construction in Section~\ref{sec:separation} gives finite dimensions
$m=330$ and $r=20N-1$ and a quadratic matrix map $Q$ satisfying
$Q(y)\succeq0$ for every $y\in\R^m$.
Proposition~\ref{prop:hyperbolic} shows that its polynomial $p_Q$ is
hyperbolic with respect to the specified direction.
Theorem~\ref{prop:no-lift} shows that the entire closed hyperbolicity cone
$K_Q$ has no finite semidefinite lift.  These are all the assertions of
Theorem~\ref{thm:main}.
\end{proof}
